% =====================================================================
%  GMN3 §1  引言与摘要（Introduction and Summary）
%  原文：Lamination-Paper.tex 360--558 行；PDF p.4--6（印刷 p.3--5）
%  编号公式 (1.1)--(1.5) 共 5 条，全部在本块内，故 \eqref 可用；
%  跨块引用（节号、式号）暂写原文编号并加 % TODO，合稿时统一改 \ref/\eqref
% =====================================================================

\section{引言与摘要}
\label{sec:IntroductionAndSummary}

本文延续一项关于模空间 $\CM$ 之性质的研究\cite{Gaiotto:2008cd,Gaiotto:2009hg}；这些模空间自然联系于四维 $\CN=2$ 超对称量子场论\cite{Seiberg:1994rs,Seiberg:1996nz}。这些模空间正是这类理论在 $\IR^3\times S^1$ 上的真空模空间。它们可以被赋予完全可积系统的结构，并且一直是大量密集研究的对象。文献 \cite{Gaiotto:2008cd,Gaiotto:2009hg} 的关注点，是 $\CM$ 上的一族函数，即所谓的「达布坐标」；这族坐标既实用又有趣。说它实用，是因为它提供了计算 $\CM$ 上\hk 度量的简洁途径，提供了可以证明 Kontsevich--Soibelman 穿墙公式\cite{ks1} 的框架，还可以用来构造强耦合下 $\CN=4$ 理论的散射振幅\cite{Alday:2009yn,Alday:2009dv}。说它有趣，是因为它具有许多漂亮的渐近性质与解析性质，把它与若干其他主题联系起来，其中包括可积系统理论的热力学 Bethe 拟设，以及（高阶）Teichm\"uller 理论的 Fock--Goncharov 坐标。

在本引言的余下部分，我们对本文的四个主要结果作一简短的说明性描述。

鉴于 $\CM$ 上达布坐标的核心重要性，为它们给出更直接的物理解释是可取的。本文的主要目标之一正是提供这样一个解释：它们是 $\CN=2$ 理论中线算子的 vev。这是我们的第一个主要结果。这一解释浓缩在第 6 节% TODO \ref{sec:linehol}
的式 (6.3)% TODO \eqref{eq:trh-large-r}
之中：
\begin{equation} \label{eq:intro-eq}
\inprod{L_\zeta} = \sum_\gamma \fro(L_\zeta, \gamma) \CY_\gamma.
\end{equation}
此式左端是一个超对称线算子 $L_\zeta$ 的真空期望值。该算子位于空间中的单个点，沿时间方向伸展，并保持由角度 $\zeta$ 指定的四个庞加莱超荷——与相位为 $\zeta$ 的 BPS 粒子所保持的正是同样的四个超对称。（第 2 节% TODO \ref{sec:Susy-BPS-line}
详细讨论这一定义，并概述超对称线算子理论的若干基本方面。）超对称 't Hooft--Wilson 算子给出了可能的 $L_\zeta$ 的例子，但正如我们将看到的，还有别的例子。式 \eqref{eq:intro-eq} 的右端是达布坐标 $\CY_\gamma$ 的展开。形象地说，$\CY_\gamma$ 可以看成低能红外理论中线算子的 vev，它由插入一个无穷重的双荷子得到；该双荷子带电磁荷 $\gamma$、中心荷相位为 $\zeta$。在这个意义上，式 \eqref{eq:intro-eq} 表达了紫外理论中的线算子 $L_\zeta$ 如何分解为红外理论中的线算子。

展开式 \eqref{eq:intro-eq} 的系数 $\fro(L_\zeta,\gamma)$ 是整数，并且实际上定义了一种新的 BPS 指标，它计数一类新的 BPS 态，我们称之为\emph{带框 BPS 态}（framed BPS state）。这些态在第 3 节% TODO \ref{sec:Hilbert-Halo-WC}
中引入。简言之，在线算子 $L$ 存在时，理论的希尔伯特空间 $\CH_L$ 是 $\N=2$ 超对称之未破缺部分的一个表示。它仍按电磁荷 $\gamma$ 分次；在标号为 $\gamma$ 的扇区中，能量满足一个修正过的 BPS 界
\begin{equation} \label{eq:bps-bound}
E\geq - {\rm Re} (Z_\gamma/\zeta).
\end{equation}
带框 BPS 态就是 $\CH_L$ 中使这个界取到等号的态。为减少混淆，我们把原 $\CN=2$ 理论中（没有线算子时）的 BPS 态称为「原味 BPS 态」（vanilla BPS state）。

我们可以精细化 BPS 指标，以计入带框 BPS 态的自旋信息。这个非常有趣的指标定义于第 3.2 节% TODO \ref{subsec:BPS-Deg-PSC}
的式 (3.12) 与式 (3.14)% TODO \eqref{eq:Def-Framed-PSC}、\eqref{eq:PSC-one}
，称为\emph{保护自旋特征标}（protected spin character，PSC）。带框 BPS 态与「原味」BPS 态都可以定义 PSC。PSC 是变量 $y$ 的函数，记作 $\fro(L_\zeta,\gamma;y)$。取 $y=-1$ 即得 BPS 指标。

本文的第二个与第三个主要结果，都围绕带框保护自旋特征标的生成函数（例如式 (3.38)% TODO \eqref{eq:form-gen}
）的性质展开：
\begin{equation}\label{eq:intro-eq-2}
F( L_\zeta ) := \sum_{\gamma} \fro( L_\zeta, \gamma;y) X_\gamma.
\end{equation}
这里生成函数中的形式变量 $X_{\gamma}$ 应视为某个非交换海森堡代数的元素，因而满足
\begin{equation} \label{eq:heis}
X_{\gamma} X_{\gamma'} = y^{\langle \gamma, \gamma' \rangle} X_{\gamma+\gamma'}
\end{equation}
其中 $\langle \gamma, \gamma' \rangle$ 是通常的 Dirac--Schwinger--Zwanziger 荷的反对称乘积。本文的第二个主要结果是：生成函数 \eqref{eq:intro-eq-2} 所成的代数，给出 $\CM$ 上函数代数的一个非交换形变，其中 $y$ 是形变参数。

这条乘法法则在第 3.6 节% TODO \ref{subsec:Deformed-Ring}
中在物理上得到论证。所得的非交换代数在第 4 节% TODO \ref{sec:Formal-line-operators}
的许多具体例子中得到展示。在第 8 节% TODO \ref{sec:Examples-of-M}
中，我们以一种便于看清「第 4 节的代数确实是第 8 节代数之非交换形变」的方式，描述达布坐标的交换环。在一类重要的例子（$A_1$ 理论）中，这个代数与量子 Teichm\"uller 理论中量子测地长度算子的代数相关，见第 11 节% TODO \ref{sec:Quantum-Holonomy}
。

PSC 与近年来被大量研究的「refined BPS 指标」\cite{Gukov:2004hz,Diaconescu:2007bf,Dimofte:2009bv,Dimofte:2009tm} 之间的主要区别在于：从\emph{先验}的理由看，PSC 是一个指标。然而在我们检查过的例子中，这两个量相符。这种相符并非平凡，它蕴涵着：无论是带框的还是「原味」的 BPS 态，都不存在超对称的某些表示。这一出人意料的相符在第 3.2 节% TODO \ref{subsec:BPS-Deg-PSC}
中被表述为一条「强正性猜想」。该猜想若成立，将产生深远影响。我们在第 4 节% TODO \ref{sec:Formal-line-operators}
中证明：强正性蕴涵线算子之环及其 vev 本质上可以用形式代数手段算出。考虑到 PSC 的几何解释，这是一个高度非平凡的结果。此外，线算子的强正性让人联想到一些深刻的数学，包括 Lusztig 与 Kashiwara 量子群的典范基、Fomin 与 Zelevinsky 簇代数的洛朗现象，以及 Fock 与 Goncharov 的泛洛朗多项式。我们并未试图把这些联系做得很精确，而是把它们留待将来的工作。

带框 BPS 态经历与「原味」BPS 态类似的穿墙。边缘稳定墙的类比物是式 (3.9)% TODO \eqref{eq:Def-BPS-Walls}
定义的 BPS 墙。在这些墙附近，一部分带框 BPS 态有一个非常简单的物理描述，它密切类比于 Frederik Denef 发现的多中心超引力解中的「halo」构型。穿墙机制与文献 \cite{Denef:2007vg} 所述本质上相同：当参数越过一道 BPS 墙时，整个 halo 构型的 Fock 空间被产生或被湮灭。本文的第三个主要结果是描述生成函数 \eqref{eq:intro-eq-2} 如何穿墙的公式。它本质上由量子双对数的共轭给出（精确公式陈述于式 (3.47) 与式 (3.48)% TODO \eqref{eq:Spin-KS-tmn-S}、\eqref{eq:Spin-KS-ii}
）。

这条带框穿墙公式的一个推论，是对 Kontsevich 与 Soibelman 关于\emph{原味} PSC 的所谓「motivic 穿墙公式」的一个简单物理推导，见第 3.5 节% TODO \ref{subsec:Motivic-WCF}
。粗略地说，想法是：通过「原味」BPS 态对带框 BPS 态的影响，来间接研究「原味」BPS 态的谱。（这与我们早先对穿墙公式的推导\cite{Gaiotto:2008cd} 有关；那个推导同样是通过「原味」BPS 态对别的东西的影响来间接研究它们——在那种情形，受影响的是维数约化理论的有效拉氏量。）这样我们就证明了：motivic KSWCF 直接来自 halo Fock 空间的简单物理。

最后，在本文余下部分，我们展示第四个主要结果：在一类有趣的 $\CN=2$ 理论例子中，如何计算带框 BPS 指标与带框 PSC；这类理论即 $\CS$ 类，由 $d=6$ 非阿贝尔 $(2,0)$ 理论在亏格 $g$、$n$ 个穿刺点的黎曼面 $C$ 上作部分扭转紧化而得到\cite{Witten:1997sc,Gaiotto:2009we,Gaiotto:2009hg}。在第 7 节% TODO \ref{sec:Six-Dim-View}
简要回顾 $d=6$ 的 $(2,0)$ 理论之后，我们利用 $d=6$ 理论的表面算子，得到使计算成为可能的关键结果 (7.13)% TODO \eqref{eq:L-5d-Holon}
。该结果断言：可以把 $\CM$ 等同于 $C$ 上平坦联络的模空间，并把线算子的量子期望值表达为平坦联络的经典和乐。在第 10 节% TODO \ref{sec:Examples}
中，我们详细展开取自「$A_1$ 理论」这一类的许多例子；这类理论是 $\CS$ 中基于 $\lieg=su(2)$ 的 $d=6$ 理论的那些。在此过程中，第 9 节% TODO \ref{sec:A1}
把共形 $A_1$ 理论中线算子的分类（由 \cite{Drukker:2009tz} 提出）推广到渐近自由理论。结论是：$C$ 上闭曲线的同痕类被所谓「叶层」（lamination）的对象取代（定义见第 9.4 节% TODO \ref{subsec:Lams-Def}
）。我们由此与 Fock 和 Goncharov 的一些非常有趣的工作\cite{MR2349682,MR2567745} 建立联系。一旦建立了这一联系，带框 BPS 指标的计算在原则上就是直接的，可以化归为一个本质上众所周知的算法——我们称之为\mbox{\emph{traffic rule}} 算法——详见附录 F% TODO \ref{app:traffic}
。

相当直接的 \mbox{traffic-rule} 计算给出了一些物理上出人意料、数学上非平凡的结果。例如，$SU(2)$、$N_f=0$ 理论中基础表示下的 Wilson 线振幅，由形如
\begin{equation} \label{eq:nf0-wilson}
\CY_{\gamma_e} + \CY_{-\gamma_e} + \CY_{\gamma_m},
\end{equation}
的三项表达式 (10.33)% TODO \eqref{eq:wilson-vev-nf0}
给出，其中 $\gamma_e$ 是 $W$ 玻色子电荷的一半，$\gamma_m$ 是一个与之互不局域的电荷。前两项是半经典推理所预期的，但第三项出人意料。它在弱耦合区域是指数小的。沿着这条思路还可以给出若干例子，见第 10 节% TODO \ref{sec:Examples}
。这里我们只再提一个例子：在强耦合的 $SU(2)$、$N_f=4$ 理论中，Wilson 线支持 $7$ 个带框 BPS 超多重态，这一点由简单的计算 (10.44)% TODO \eqref{eq:Nf4-Wilson}
揭示。要通过研究单极子模空间的几何来重现这个结果，大概会相当不平凡。

在第 5 节与第 12 节% TODO \ref{sec:Cluster-Algebras}、\ref{sec:Tropic}
中，我们简要探讨与一些有趣数学的联系，并说明这些联系有实用的物理应用。在第 5 节% TODO \ref{sec:Cluster-Algebras}
中我们说明，形式线算子的代数是一个量子簇代数的例子。这一观点在构造 $\CN=2^*$ 理论的形式线算子时很有用，并对簇代数文献中一些正性猜想提出了精细化。我们引入「簇 $\CN=2$ 理论」的概念。在第 12 节% TODO \ref{sec:Tropic}
中，我们类似地引入「热带 $\CN=2$ 理论」的概念，并为线算子引入「热带标签」。应当把它们看成是对（单）线算子之枚举的一种紫外到红外的映射。在这两节中，我们只是触及了也许非常深刻的联系的表层。

一些技术性论证与一些约定被放进附录。

尽管篇幅（令人遗憾地）很长，本文仍留下许多漂亮的石头没有翻动。我们在第 13 节% TODO \ref{sec:Open-Problems}
指出若干开放问题与未来的研究方向。
